All
Start with a template

\documentclass[a4paper,11pt,oneside]{book}

% ==========================================
% PACKAGES
% ==========================================

\usepackage[T1]{fontenc}
\usepackage[english]{babel}

\usepackage{graphicx}
\usepackage{amsmath}
\usepackage{amssymb}
\usepackage{float}
\usepackage{booktabs}
\usepackage{xcolor}
\usepackage{xurl}
\usepackage{hyperref}

\graphicspath{{figures/}}
\setlength{\emergencystretch}{3em}

% ==========================================
% DOCUMENT INFORMATION
% ==========================================

\title{VISUAL-BASED CONTROL OF AN AGILEX LIMO}
\author{Giulio Cataldo}
\date{\today}

\begin{document}

\frontmatter

\hypersetup{pageanchor=false}
\maketitle
\hypersetup{pageanchor=true}
\tableofcontents

\mainmatter

% ==========================================
% CHAPTER 1 - INTRODUCTION
% ==========================================

\chapter{Introduction}

In May 2026, I started a six-month internship at CRAN (Centre de Recherche en Automatique de Nancy), a research laboratory affiliated with the University of Lorraine in Nancy, France.

The project consisted of two main tasks. Given an AgileX LIMO mobile robot and a scaled-down urban environment, my first objective was to set up a suitable software environment to facilitate future students' work on the platform.

The second and more challenging objective was to develop a complete planning and control pipeline capable of autonomously navigating the robot from a starting point to a designated destination, while avoiding physical obstacles and complying with basic traffic rules.

The main outcomes achieved at the end of the internship are summarized as follows:

\begin{enumerate}
    \item \textbf{Gazebo Simulation Environment:}
    Development of a high-fidelity Gazebo simulation framework featuring a digital model of the LIMO robot and a virtual replica of the laboratory's test track.

    \item \textbf{Autonomous Navigation Pipeline:}
    Implementation of an end-to-end software architecture integrating global mapping, sensor-fused localization, trajectory generation, and optimal control, allowing the robot to reach designated waypoints while avoiding both physical barriers and visual road obstacles.

    \item \textbf{System Documentation and Deployment:}
    Compilation of a detailed, step-by-step setup guide to streamline the installation of all necessary ROS 2 packages, sensor configurations, and system launch procedures via custom scripts and launch parameters.
\end{enumerate}

% ==========================================
% CHAPTER 2 - SYSTEM OVERVIEW
% ==========================================

\chapter{System Overview}

\section{The AgileX LIMO Robot}
\begin{figure}[H]
    \centering
    \includegraphics[width=0.5\linewidth]{limo.png}
    \caption{The Limo Robot}
    \label{fig:limo_robot}
\end{figure}
\subsection{Vehicle Specifications}

The AgileX LIMO is a compact mobile robotic platform designed
for educational and research applications. Its main technical
specifications are summarized in Table~\ref{tab:limo_specs}.



\begin{table}[H]
    \centering
    \caption{Main technical specifications of the AgileX LIMO.}
    \label{tab:limo_specs}

    \begin{tabular}{ll}
        \toprule
        \textbf{Parameter} & \textbf{Specification} \\
        \midrule
        Weight & 4.2 kg \\
        Dimensions & $322 \times 215 \times 247$ mm \\
        Run time & 40 minutes \\
        Standby time & 2 hours \\
        Operating system & Ubuntu 18.04 \\
        Maximum climbable slope & $25^\circ$ \\
        Operating temperature & $-10^\circ$C to $40^\circ$C \\
        Charging interface & 5.5 $\times$ 2.1 mm DC barrel jack \\
        CPU & ARM64 Quad-Core 1.43 GHz (Cortex-A57) \\
        GPU & 128-core NVIDIA Maxwell (921 MHz) \\
        Battery & 12 V Li-ion, 5600 mAh \\
        \bottomrule
    \end{tabular}
\end{table}

The AgileX LIMO platform supports multiple locomotion modalities;
for this project, the \textbf{Ackermann steering mode} was selected
to closely replicate standard automotive kinematics.
Under this configuration, the vehicle features a minimum turning
radius of $0.4\,\mathrm{m}$.

\subsection{Onboard Computing Hardware}

The robot's primary compute unit is an \textbf{NVIDIA Jetson Nano (4GB)}.
It is worth noting that the onboard CPU exhibits limited computational
capabilities, making it a critical constraint for resource-intensive
processes. While the integrated GPU can support hardware-accelerated
tasks---such as lightweight neural network inference or visualization
rendering---it was not directly utilized within the scope of the
developed software pipeline, leaving the execution to rely entirely
on CPU-side optimizations.

\subsection{Software Architecture}

The original software environment used to operate the robot and
implement control algorithms was based on a hybrid architecture
combining \textbf{ROS 1 Melodic} and \textbf{ROS 2 Eloquent}.
Communication between the two distributions was established through
the \texttt{ros1\_bridge} package, which enabled the exchange of
messages across their respective topics.

The robot description files and primary launch configurations were
provided by the official AgileX repository, originally
developed for ROS 1. Consequently, the robot's initialization
and hardware configuration had to be performed within the ROS 1
environment, while additional nodes could be developed in ROS 2
and connected to the existing system through the bridge.

However, this architecture presented significant limitations,
primarily due to the communication overhead introduced by the bridge
and the outdated software distributions.

In particular, the transmission of large messages, such as camera
images, resulted in reduced message frequency and increased temporal
jitter on the ROS 2 side. These limitations made the architecture
unsuitable for the reliable execution of real-time computer vision
algorithms.

Furthermore, the \texttt{ros1\_bridge} process consumed a
considerable portion of the Jetson Nano's already limited CPU
resources, significantly restricting the computational capacity
available for more advanced navigation and control algorithms.


(source: \url{https://github.com/agilexrobotics/limo-doc/blob/master/Limo\%20user\%20manual(EN).md})
(source: \url{https://github.com/agilexrobotics/limo_ros})
(source: \url{https://docs.trossenrobotics.com/agilex_limo_docs/specifications/limo.html})

\section{Sensor Overview}

The AgileX LIMO is equipped with several sensors that provide
information about the robot's motion and its surrounding
environment. The main sensors employed in this project are
an inertial measurement unit (IMU), a 2D LiDAR, and an
RGB-D camera.

\subsection{Inertial Measurement Unit (IMU)}

The robot is equipped with an \textbf{MPU6050 IMU}, which
integrates a three-axis accelerometer and a three-axis
gyroscope.

The accelerometer measures the specific force acting on
the robot, while the gyroscope provides its angular
velocity.

In the proposed architecture, the IMU measurements are
combined with wheel odometry through an Extended Kalman
Filter (EKF) to improve the estimation of the robot's
motion.

\subsection{LiDAR}

The robot is equipped with an \textbf{EAI X2L 2D LiDAR},
which measures the distance to surrounding objects
by scanning the environment in a horizontal plane.

\begin{table}[H]
    \centering
    \caption{Main technical specifications of the EAI X2L LiDAR.}
    \label{tab:lidar_specs}
    
    \begin{tabular}{lcccc}
        \toprule
        \textbf{Parameter} & \textbf{Minimum} & \textbf{Default} & \textbf{Maximum} & \textbf{Unit} \\
        \midrule
        Ranging frequency & -- & 3000 & -- & Hz \\
        Scanning frequency & 5 & 6 & 8 & Hz \\
        Ranging range & 0.12 & -- & 8 & m \\
        Scanning angle & -- & 0--360 & -- & $^\circ$ \\
        Absolute tolerance & -- & 2 & -- & cm \\
        Relative tolerance & -- & 3.5 & -- & \% \\
        Pitch angle & 0.25 & 1 & 1.75 & $^\circ$ \\
        Angular resolution & 0.60 (5 Hz) & 0.72 (6 Hz) & 0.96 (8 Hz) & $^\circ$ \\
        \bottomrule
    \end{tabular}
\end{table}

\subsection{RGB-D Camera}

The \textbf{Orbbec DaBai Stereo Depth Camera} provides
both color images and depth measurements of the
surrounding environment.

The camera employs structured-light technology to
estimate depth by projecting an infrared pattern
onto the environment and analyzing its deformation
in the acquired images.

However, this sensing technology presents several
limitations. In particular, surfaces with low
infrared reflectivity, strong ambient illumination,
and interference from other structured-light cameras
may result in missing or unreliable depth
measurements.

These limitations must be considered when using
depth information for obstacle detection and
environment mapping.

\begin{table}[H]
    \centering
    \caption{Main technical specifications of the Orbbec DaBai Stereo Depth Camera.}
    \label{tab:camera_specs}

    \begin{tabular}{ll}
        \toprule
        \textbf{Parameter} & \textbf{Specification} \\
        \midrule
        Infrared camera baseline & 40 mm \\
        Depth range & 0.3--3 m \\
        Depth resolution & $640 \times 400$ @ 30 FPS \\
                         & $320 \times 200$ @ 30 FPS \\
        RGB resolution & $1920 \times 1080$ @ 30 FPS \\
                       & $1280 \times 720$ @ 30 FPS \\
                       & $640 \times 480$ @ 30 FPS \\
        Depth accuracy & 6 mm at 1 m \\
        Depth horizontal FOV & $67.9^\circ$ \\
        Depth vertical FOV & $53^\circ$ \\
        RGB horizontal FOV & $71^\circ$ \\
        RGB vertical FOV & $43.7^\circ$ \\
        Latency & 30--45 ms \\
        Data transmission & USB 2.0 or higher \\
        \bottomrule
    \end{tabular}
\end{table}

(source: \url{https://www.digikey.fr/fr/products/detail/tdk-invensense/MPU-6050/4038009?srsltid=AYCX9vLdGtob9dv5gB5Dq7AO9A5o5Ce9G7mWNG-w6T4ix5jYLCnoezjk})
(source: \url{https://virtuale.unibo.it/pluginfile.php/2664001/mod_resource/content/3/sensors_overview.pdf})
(source: \url{https://robotics.growbotics.ai/suppliers/products/orbbec-dabai-depth-camera})

\section{Laboratory Circuit}

The robot is designed to navigate within a laboratory test track
developed by Quanser, featuring approximate dimensions of
$5 \times 6\,\mathrm{m}$.

\begin{figure}[H]
    \centering
    \includegraphics[width=0.8\linewidth]{circuit.png}
    \caption{Plan of the circuit}
    \label{fig:laboratory_circuit_plan}
\end{figure}

Since the track is entirely flat, compliance with traffic regulations strictly requires the integration of computer vision capabilities. Furthermore, the platform perimeter is bounded by walls on only three of its four sides. 

This partial enclosure presents both advantages and technical challenges for the localization subsystem:
\begin{itemize}
    \item \textbf{Geometric anchors:} Irregularities along the perimeter provide unique features that enhance LiDAR-based localization accuracy.
    \item \textbf{Environmental noise:} Conversely, the open sections of the circuit introduce inconsistent features into the spatial map---as detailed in the section on map creation---rendering the localization process inherently unstable.
\end{itemize}

To mitigate these instability issues and prevent mapping artifacts, uniform white obstacles were positioned to mask the critical outer corners.

\begin{figure}[H]
    \centering
    \includegraphics[width=0.5\linewidth]{blank.png}
    \caption{Real circuit environment}
    \label{fig:real_circuit_environment}
\end{figure}

The circuit presents in addition a very reflective surface that, combined with the lack of proper coverage from sunlight and poor dynamic range of the camera, yields either relevant lighting variations or directly to partial blindness of the vision.

% ==========================================
% CHAPTER 3 - ENVIRONMENT SETUP
% ==========================================

\chapter{Environment Setup}
In this section we will explain the changes actuated for the optimization of the robot's software, the configuration of the sensors and the creation of an independent simulation environment.

\section{Robot Software Upgrade}
The first challenge addressed was the presence of the \texttt{ros1\_bridge}. Operating exclusively within ROS 1 presented significant limitations: as an outdated framework, relying solely on this stack did not align with personal skill development goals and, above all, proved counterproductive to the educational and technical objectives of the project. 

On the other hand, transitioning directly to ROS 2 Eloquent was unfeasible, as the lack of supported package repositories to initialize the system for that release made a comprehensive system upgrade necessary.

Fortunately, a practical solution was found in a custom Ubuntu 20.04 system image developed and publicly released by Qengineering. This provided a reliable path for upgrading the Jetson Nano's operating system while fully preserving compatibility with the NVIDIA software stack and drivers. 

This upgrade enabled a complete migration from the hybrid ROS 1--ROS 2 architecture to \textbf{ROS 2 Foxy}, a distribution that shares core architectural principles and key components with more recent LTS releases such as ROS 2 Humble.

\section{Robot and Sensor Bringup}

\subsection{Sensor Configuration}

To properly bring up the robot and its sensors, the official AgileX ROS 2 repository was used as a starting point.

(source: \url{https://github.com/agilexrobotics/limo_ros2/tree/foxy})

Additionally, the software packages provided by the sensor manufacturers were used to establish communication with the LiDAR and the RGB-D camera.

(source: \url{https://github.com/orbbec/ros2_astra_camera}, \url{https://github.com/orbbec/ros2_astra_camera/tree/master/astra_camera/openni2_redist/arm64})

Despite the use of official drivers, several configuration adjustments and minor source code modifications remained necessary to ensure the proper operation of the robotic platform and its sensor suite. Furthermore, the source code was refactored to make topic names fully configurable, thereby guaranteeing greater consistency and seamless interoperability between the physical robot and the simulation environment.

\subsection{Odometry and IMU Fusion}

Another critical enhancement implemented is the sensor fusion pipeline, which combines wheel odometry and IMU measurements via an Extended Kalman Filter (EKF) provided by the official \texttt{robot\_localization} package. The following boolean table shows if the contribution of the sensor is enabled for the estimation of the kinematic quantities.

\begin{table}[H]
    \centering
    \begin{tabular}{lcc}
        \toprule
        \textbf{Quantity} & \textbf{Odom (\texttt{/odom})}
                         & \textbf{IMU (\texttt{/limo/imu})} \\
        \midrule
        Position $x, y$                       & Yes & No \\
        Yaw angle $\psi$                      & Yes & No \\
        Linear velocities $v_x, v_y$          & Yes & No \\
        Yaw rate $\dot{\psi}$                 & No  & Yes\\
        Height $z$, roll and pitch            & No  & No \\
        Vertical velocity                    & No  & No \\
        Roll and pitch rates                 & No  & No \\
        Accelerations $a_x, a_y, a_z$          & No  & No \\
        \bottomrule
    \end{tabular}
    \caption{Quantities used from the odometry and IMU messages.}
    \label{tab:odom_imu}
\end{table}

\section{Simulation Setup}

\subsection{Robot model completion}

The main challenge of working on a ROS~2 Foxy platform was the lack
of suitable robot description files for launching a simulation
environment.

To solve this problem, the \texttt{limo\_car} folder and the
\path{limo_start.launch.py} launch file were imported from the
Humble branch of the LIMO repository and adapted to the Foxy
workspace.

\subsection{Ackermann model compatibility}

Furthermore, an auxiliary ROS~2 node,
\texttt{gazebo\_twist\_adapter.py}, was introduced to ensure
compatibility between the velocity commands generated by
\texttt{teleop\_twist\_keyboard} and the Ackermann steering plugin
used in Gazebo.

While the teleoperation node publishes linear and angular velocity
commands, the Ackermann plugin requires longitudinal velocity and
steering angle as inputs. Therefore, an intermediate conversion
is necessary.

The adapter performs this conversion using the kinematic bicycle
model.

\begin{figure}[H]
    \centering
    \includegraphics[width=0.85\textwidth]{blank.png}
    \caption{Command conversion node and associated ROS topics.}
    \label{fig:twist_adapter}
\end{figure}

\begin{figure}[H]
    \centering
    \includegraphics[width=0.75\textwidth]{blank.png}
    \caption{Ackermann steering configuration and equivalent
    kinematic bicycle model.}
    \label{fig:ackermann_bicycle}
\end{figure}

Figure~\ref{fig:ackermann_bicycle} illustrates the relationship
between the Ackermann steering configuration and the kinematic
bicycle model. In this representation, $\omega$ denotes the yaw
rate, $v$ the longitudinal velocity, $\delta$ the equivalent
front-wheel steering angle, and $L$ the wheelbase, i.e., the
distance between the front and rear axles.

The relationship between the yaw rate and the steering angle is
\begin{equation}
    \dot{\theta} = \omega = \frac{v}{L}\tan\delta.
    \label{eq:bicycle_yaw_rate}
\end{equation}

Rearranging Equation~\eqref{eq:bicycle_yaw_rate} yields the steering
angle required by the Ackermann plugin:
\begin{equation}
    \delta_{\mathrm{plugin}}
    = \arctan\left(\frac{L\omega}{v}\right).
    \label{eq:ackermann_steering_angle}
\end{equation}

source:
\begin{quote}
    \url{https://virtuale.unibo.it/pluginfile.php/2494729/mod_resource/content/1/MoMfA_1_Introduction.pdf}
\end{quote}

\subsection{Gazebo Environment}

The simulated robot operates within a
\textbf{virtual road network} with approximate dimensions of
$12 \times 7\,\mathrm{m}$, manually designed using Inkscape to
reproduce the main features of the experimental track available
in the laboratory. For simplicity, the simulated environment is
enclosed by uniformly shaped walls, providing a simplified
representation of the physical boundaries.

\begin{figure}[H]
    \centering
    \includegraphics[width=0.85\textwidth]{simulation_circuit_no_obst.png}
    \caption{Simulated road network in Gazebo.}
    \label{fig:gazebo_track}
\end{figure}

In addition to this Gazebo environment, a second version was
developed by introducing a few physical obstacles along the track.
This alternative configuration was designed to evaluate the
navigation system's ability to simultaneously handle physical
obstacles detected by the LiDAR and road markings identified
through computer vision.

\begin{figure}[H]
    \centering
    \includegraphics[width=0.85\textwidth]{simulation_circuit_obst.png}
    \caption{Simulated road network with additional physical obstacles.}
    \label{fig:gazebo_track_obstacles}
\end{figure}

\subsection{Extended Kalman Filter in Simulation}

The Extended Kalman Filter (EKF) was also implemented in the simulation environment to maintain consistency with the real robot's localization architecture. However, its introduction did not lead to significant improvements in odometry estimation, as Gazebo provides idealized odometry measurements under the adopted simulation configuration.

\section{Environment Validation}

Thanks to these changes, both the real robot and the simulation can start up properly and, most importantly, by aligning the names of the two environments, the same code can now be used in both simulation and on the physical robot, provided that computational limitations are not taken into account.

\begin{figure}[H]
    \centering
    \includegraphics[width=0.85\textwidth]{blank.png}
    \caption{Grafo dei topic e dei frame in simulazione e nella realtà.}
    \label{fig:environment_topics_frames}
\end{figure}

\subsection{Workspace Preparation for Laboratory Courses}

The procedure described above can then be applied to each robot in the laboratory, providing students with a consistent platform for experimentation. To support this process, all required files have been uploaded to the University of Lorraine's GitLab instance, together with shell scripts (\texttt{.sh}) that automate dependency installation, sensor configuration, and the creation of convenient command aliases.

% ==========================================
% PART I - SIMULATION
% ==========================================

\part{Simulation}

% ==========================================
% CHAPTER 4 - AUTONOMOUS NAVIGATION
% ==========================================

\chapter{Simulation Pipeline}
\label{chap:simulation_pipeline}

This chapter presents the development of a complete simulation
pipeline comprising mapping, localization, path planning, and
control. The pipeline operates in the simulated environments
introduced previously, using data from the simulated camera
and LiDAR.

\section{General Control Architecture}

The general control architecture of the mobile robot is shown
in Figure~\ref{fig:general_control_scheme}.

\begin{figure}[H]
    \centering
    \includegraphics[width=\textwidth]{gen_ctrl_scheme.png}
    \caption{General control architecture of the mobile robot.}
    \label{fig:general_control_scheme}
\end{figure}

The user assigns a task to the robot, such as reaching a desired
target pose. Based on this objective and the available sensor
information, the planner generates a reference path for the robot
to follow. The controller then computes the appropriate control
commands based on the reference path and the estimated robot state.

These commands are sent to the robot's actuators, producing changes
in its position and orientation. The resulting motion is measured
by the onboard sensors, allowing the robot's state estimate and
control commands to be continuously updated. This feedback loop
continues until the assigned task is completed.

\section{Implemented Navigation Pipeline}

The navigation pipeline implemented in simulation is illustrated
in Figure~\ref{fig:simulation_pipeline}.

\begin{figure}[H]
    \centering
    \includegraphics[width=\textwidth]{blank.png}
    \caption{Overview of the navigation pipeline implemented
    in simulation.}
    \label{fig:simulation_pipeline}
\end{figure}

Semantic labels extracted from camera images are combined with
depth data to generate a three-dimensional semantic point cloud,
in which each point carries both spatial coordinates and a semantic
class. During mapping, these observations are integrated with the
LiDAR-based occupancy map to construct a global representation of
the environment, capturing both obstacle geometry and semantic
traversability information.

Navigation begins when the user specifies a target pose, comprising
position and orientation. The global planner then computes a
reference path using \textbf{Smac Hybrid A*} with a
\textbf{Dubins motion model}, accounting for the robot's Ackermann
steering geometry and minimum turning radius. During navigation,
incoming semantic observations and LiDAR scans continuously update
a local costmap, representing nearby obstacles and changes in
the environment.

The reference path, local costmap, and estimated robot state are
provided to a \textbf{sampling-based nonlinear Model Predictive
Controller (MPC)}. At each control cycle, the controller evaluates
candidate sequences of velocities and steering angles by predicting
the resulting robot motion over a finite horizon. It selects the
best admissible sequence among the evaluated candidates according
to path-tracking, obstacle-avoidance, and control-effort costs,
and applies its first reachable command.

This process repeats until the robot reaches the target within
the specified position and orientation tolerances, or navigation
is interrupted by an error or a stop request.

The following sections describe and justify each component of
the pipeline, presenting its inputs, outputs, and corresponding
experimental results.

\section{Image Processing}
\label{sec:computer_vision_pipeline}

The main objective of the project, as previously mentioned, is to
enable the robot to navigate the experimental track while complying
with a simplified set of traffic rules. Achieving this objective
requires the robot to recognize and interpret road markings, making
computer vision a fundamental component of the navigation system.

Furthermore, since the experimental track is assumed to be flat,
the geometric relationship between camera observations and the
ground plane can be exploited to simplify the spatial reconstruction
of the detected road markings.

The image-processing system was implemented within the ROS~2 package
\texttt{cv\_package}. It consists of several interconnected nodes
forming a dedicated perception pipeline that processes RGB-D images
and converts the extracted semantic information into a representation
compatible with the LiDAR-based navigation system.

The overall architecture of this subsystem is illustrated in
Figure~\ref{fig:cv_pipeline}.

\begin{figure}[H]
    \centering
    \includegraphics[width=\textwidth]{blank.png}
    \caption{Architecture of the computer vision pipeline
    implemented in \texttt{cv\_package}.}
    \label{fig:cv_pipeline}
\end{figure}

The following sections describe each processing node, focusing on
its operating principles, input and output data, and implementation
choices.

\subsection{Image Thresholding: \texttt{lane\_detector.py}}
\label{subsec:lane_detector}

Four main semantic elements can be distinguished in both the
simulated and real tracks: the dark grey road surface, visual
obstacles (typically white or grey), yellow lines, and white dashed
lines. Due to the robot's limited onboard computational resources,
a lightweight colour-based detection algorithm was selected
instead of a deep neural network.

The adopted strategy first isolates the road surface and yellow
lines. These features exhibit distinctive and relatively consistent
colour profiles in both environments, making them easier to
identify than the remaining elements, which are extracted in
subsequent processing stages.

The purpose of this node is therefore to process incoming camera
frames and generate semantic masks for dark grey and yellow regions
with minimal computational overhead.

\subsubsection{Input and Preprocessing}

The node receives RGB images with a resolution of
$320 \times 240$ pixels at a frequency of $20\,\mathrm{Hz}$ through
the \texttt{/rgb/image\_raw} topic.

The upper half of each image is
discarded, as it contains no relevant road information. Depending
on the camera, a configurable number of bottom rows may also be
excluded, since these can contain consistently black pixels.
These operations reduce the computational load of subsequent
processing stages and prevent image artefacts from affecting
the segmentation.

The retained image region is then converted from RGB to HSV.
Separating hue, saturation, and brightness facilitates the
identification of yellow regions and allows dark pixels to be
selected according to their brightness. 

\subsubsection{Colour Segmentation}

The HSV thresholds used for segmentation are reported in
Table~\ref{tab:lane_detector_thresholds}. Hue values follow the
OpenCV convention, ranging from 0 to 179, while saturation and
value range from 0 to 255.

\begin{table}[H]
    \centering
    \begin{tabular}{lccc}
        \toprule
        \textbf{Class} & \textbf{H} & \textbf{S} & \textbf{V} \\
        \midrule
        Dark pixels & Any & Any & $0$--$150$ \\
        Yellow      & $15$--$35$ & $80$--$255$ & $80$--$255$ \\
        \bottomrule
    \end{tabular}
    \caption{HSV thresholds used by \texttt{lane\_detector.py}.}
    \label{tab:lane_detector_thresholds}
\end{table}

Dark pixels are classified as road candidates, while pixels within
the yellow colour range are identified as yellow markings.
All remaining valid pixels are assigned to the background class.

\subsubsection{Output Representation}

The segmentation results are encoded in a single-channel grayscale
image, with each semantic label represented by a distinct intensity
value. The output classes are listed in
Table~\ref{tab:lane_detector_labels}.

\begin{table}[H]
    \centering
    \begin{tabular}{lp{0.53\textwidth}}
        \toprule
        \textbf{Label} & \textbf{Description} \\
        \midrule
        \texttt{LABEL\_INVALID}
        & Ignored pixels in the excluded bottom rows. \\

        \texttt{LABEL\_ROAD}
        & Dark pixels identified as road. \\

        \texttt{LABEL\_YELLOW}
        & Pixels identified as yellow markings. \\

        \texttt{LABEL\_BACKGROUND}
        & All other valid pixels. \\
        \bottomrule
    \end{tabular}
    \caption{Semantic labels produced by \texttt{lane\_detector.py}.}
    \label{tab:lane_detector_labels}
\end{table}

\subsection{Depth Data Processing: \texttt{Depth\_correction.py}}
\label{subsec:depth_correction}

This helper node adapts the depth images to the dimensions and
cropping applied to the semantic images produced by
\texttt{lane\_detector.py}.

The input depth image is cropped and downsampled to obtain the
required output resolution. No interpolation is performed.
Instead, the minimum depth within each downsampling region is
retained, preserving the closest observed surface.

% TODO: Validate and describe the exact cropping, downsampling,
% treatment of invalid depth values, and pixel correspondence
% with the semantic images.

\subsection{Spatial Representation: \texttt{visual\_ptcld.py}}
\label{subsec:visual_ptcld}

The semantic images identify the road surface and yellow markings,
but do not yet distinguish white dashed lines from visually similar
obstacles. Moreover, the image coordinates alone do not directly
provide the metric spatial information required for mapping.

The purpose of \texttt{visual\_ptcld.py} is therefore to combine
semantic labels with depth measurements to reconstruct spatial
points that can be integrated with LiDAR observations for global
and local mapping. The resulting spatial representation is also
used to distinguish the boardwalk from white dashed lines.

\subsubsection{Reprojection and Ground-Plane Approximation}

The procedure first separates interior road pixels from pixels
along the boundary between the road and background classes.
The selected pixels are then downsampled and associated with
the corresponding measurements in the aligned depth image.

Using the intrinsics of the camera, these observations are reprojected
into three-dimensional points in the camera frame. The points
are subsequently transformed into the robot's base frame.

Assuming that the relevant visual features lie approximately
on the ground plane, the vertical coordinate of each transformed
point is set to zero:
\begin{equation}
    \mathbf{p}_{\mathrm{base}}
    =
    \begin{bmatrix}
        x & y & z
    \end{bmatrix}^{\mathsf{T}}
    \quad \longrightarrow \quad
    \mathbf{p}_{\mathrm{BEV}}
    =
    \begin{bmatrix}
        x & y & 0
    \end{bmatrix}^{\mathsf{T}}.
    \label{eq:ground_plane_projection}
\end{equation}

This approximation produces a planar bird's-eye view (BEV)
representation, simplifying subsequent spatial processing
and mapping. Downsampling further reduces the computational
load and memory requirements.

\subsubsection{Spatial Classification}

Once the points have been projected onto the ground plane,
spatial proximity is used to distinguish boardwalk regions
from white dashed lines.

Let $\mathcal{R}$ denote the set of road-boundary points and
$\mathcal{B}$ the set of background points. The classification
procedure comprises the following steps:

\begin{enumerate}
    \item For each background point $\mathbf{b}\in\mathcal{B}$,
    compute its Euclidean distance to the nearest road-boundary
    point:
    \begin{equation}
        d_{\mathcal{R}}(\mathbf{b})
        =
        \min_{\mathbf{r}\in\mathcal{R}}
        \left\|\mathbf{b}-\mathbf{r}\right\|_2.
    \end{equation}

    \item Classify background points whose distance exceeds
    a threshold $d_1$ as boardwalk, forming an initial set
    $\mathcal{T}$:
    \begin{equation}
        \mathcal{T}
        =
        \left\{
        \mathbf{b}\in\mathcal{B}
        \;:\;
        d_{\mathcal{R}}(\mathbf{b})>d_1
        \right\}.
    \end{equation}

    \item For each remaining background point, compute the
    Euclidean distance to the nearest point in the initial
    boardwalk set:
    \begin{equation}
        d_{\mathcal{T}}(\mathbf{b})
        =
        \min_{\mathbf{t}\in\mathcal{T}}
        \left\|\mathbf{b}-\mathbf{t}\right\|_2.
    \end{equation}

    \item Extend the boardwalk classification to background
    points whose distance to $\mathcal{T}$ is below a second
    threshold $d_2$.
\end{enumerate}

KD-Trees are used to accelerate the nearest-neighbour searches.
Writing $R=|\mathcal{R}|$, $B=|\mathcal{B}|$, and
$T=|\mathcal{T}|$, the indicative average-case costs are
$\mathcal{O}(R\log R)$ for constructing the road-boundary tree
and $\mathcal{O}(B\log R)$ for querying the background points.
The second stage requires constructing a tree from the
boardwalk points, with a cost of $\mathcal{O}(T\log T)$,
followed by at most $B$ queries, with an indicative cost of
$\mathcal{O}(B\log T)$. These query estimates depend on the
point distribution and are not worst-case guarantees.

% TODO: Verify the complexity assumptions against the selected
% reference and the KD-Tree implementation before final submission.
% Reference supplied:
% https://arxiv.org/pdf/cs/9901013

\subsubsection{Output Representation}

The resulting semantic classes are summarized in
Table~\ref{tab:visual_ptcld_labels}.

\begin{table}[H]
    \centering
    \small
    \begin{tabular}{lp{0.48\textwidth}}
        \toprule
        \textbf{Label} & \textbf{Description} \\
        \midrule
        \texttt{LABEL\_INVALID}
        & Invalid or ignored observations. \\

        \texttt{LABEL\_ROAD}
        & Road-boundary points. \\

        \texttt{LABEL\_YELLOW}
        & Yellow marking points. \\

        \texttt{LABEL\_BACKGROUND}
        & Remaining background points, interpreted as
        white dashed lines. \\

        \texttt{LABEL\_INTERIOR\_ROAD}
        & Interior road points. \\

        \texttt{LABEL\_BOARDWALK}
        & Boardwalk points within a threshold distance
        of the road. \\

        \texttt{LABEL\_INTERIOR\_BOARDWALK}
        & Remaining boardwalk points. \\
        \bottomrule
    \end{tabular}
    \caption{Semantic labels produced by
    \texttt{visual\_ptcld.py}.}
    \label{tab:visual_ptcld_labels}
\end{table}

Following classification and further downsampling, the points
are published in a \texttt{sensor\_msgs/msg/PointCloud2} message.
This representation supports visual inspection in RViz2 and
preserves the distinction between observed free space,
represented by road points, and unobserved regions, where
no valid points are available. In particular, gaps in the
projected observations are not automatically interpreted
as free space.

\begin{figure}[H]
    \centering
    \includegraphics[width=\textwidth]{blank.png}
    \caption{Semantic point cloud generated by
    \texttt{visual\_ptcld.py}.}
    \label{fig:semantic_point_cloud}
\end{figure}

Figure~\ref{fig:cv_data_flow} summarizes the input
and output topics, image dimensions, and publishing frequencies of the computer vision pipeline when operating in standalone mode.

\begin{figure}[H]
    \centering
    \includegraphics[width=\textwidth]{blank.png}
    \caption{Input and output data flow of CV pipeline, including image dimensions and
    publishing frequencies in standalone operation.}
    \label{fig:cv_data_flow}
\end{figure}

% TODO: Add measured processing times and output frequencies
% after experimental validation.

\section{Offline Global Mapping}
\label{sec:offline_mapping}

The purpose of this stage is to combine the semantic point cloud
generated by the perception pipeline with LiDAR observations to
construct a global costmap of the circuit.

\subsection{Visual Cost Map: \texttt{semantic\_grid.py}}
\label{subsec:semantic_grid}

The semantic point cloud contains several classes, which are
associated with different traversal costs. Using a cost range
of $[0,100]$, the assignments are summarized in
Table~\ref{tab:semantic_costs}.

\begin{table}[H]
    \centering
    \begin{tabular}{lc}
        \toprule
        \textbf{Semantic label} & \textbf{Cost} \\
        \midrule
        \texttt{LABEL\_INVALID}            & Discarded \\
        \texttt{LABEL\_ROAD}               & 0 \\
        \texttt{LABEL\_INTERIOR\_ROAD}      & 0 \\
        \texttt{LABEL\_BACKGROUND}         & 30 \\
        \texttt{LABEL\_YELLOW}             & 60 \\
        \texttt{LABEL\_BOARDWALK}          & 90 \\
        \texttt{LABEL\_INTERIOR\_BOARDWALK} & 90 \\
        \bottomrule
    \end{tabular}
    \caption{Traversal costs assigned to the semantic classes.}
    \label{tab:semantic_costs}
\end{table}

The global map is discretized into cells measuring
$5 \times 5\,\mathrm{cm}$, each of which is assigned a single
semantic cost.

\subsubsection{Temporal Evidence Accumulation}

A single cell may receive observations belonging to different
semantic classes, both within the same point cloud and over time.
To combine these observations, a temporal evidence accumulator
was implemented, taking inspiration from the Bayesian log-odds
filter described in \textit{Probabilistic Robotics},
Section~4.1.4.

For each cell, a separate score is maintained for each candidate
cost label. Whenever new visual observations are available,
the score associated with class $c$ is updated as follows:
\begin{equation}
    s_{t,c}
    =
    \operatorname{clip}
    \left(
        s_{t-1,c}
        + h f_{t,c}
        - m\left(1-f_{t,c}\right),
        -L,
        L
    \right),
    \label{eq:semantic_score_update}
\end{equation}
where
\begin{equation}
    f_{t,c}
    =
    \frac{N_{t,c}}{N_t}.
    \label{eq:semantic_class_fraction}
\end{equation}

Here, $N_{t,c}$ is the number of points assigned to class $c$
within the cell at time $t$, and $N_t$ is the total number of
valid points observed in that cell. The parameters $h$ and $m$
control the positive and negative score increments, respectively,
while $L$ limits the score magnitude.

The clipping operation is defined as
\begin{equation}
    \operatorname{clip}(a,-L,L)
    =
    \min\left(L,\max\left(-L,a\right)\right).
\end{equation}

The cell is assigned the class with the highest score, provided
that this score exceeds the classification threshold. Otherwise,
the cell remains unknown.

Score saturation prevents evidence from accumulating without
bound, allowing the classifier to respond to changes even after
a prolonged period of consistent observations.

The reference used as inspiration for this accumulation strategy
is available at:
\begin{quote}
    \url{https://docs.ufpr.br/~danielsantos/ProbabilisticRobotics.pdf}
\end{quote}

\subsubsection{Experimental Evaluation}

To evaluate the effectiveness of the filter, a ROS bag containing
the topics required for mapping was recorded. This recording
allows different parameter configurations to be compared using
the same input observations.

\begin{figure}[H]
    \centering
    \includegraphics[width=\textwidth]{blank.png}
    \caption{Comparison of visual maps obtained using different
    temporal evidence accumulation parameters.}
    \label{fig:semantic_grid_parameter_comparison}
\end{figure}

\subsection{LiDAR Mapping}
\label{subsec:lidar_mapping}

LiDAR observations serve two main purposes during mapping:
\begin{enumerate}
    \item Correcting the estimated robot trajectory through
    Simultaneous Localization and Mapping (SLAM).
    \item Recording physical obstacles in the costmap.
\end{enumerate}

The robot trajectory estimated from odometry is corrected online
using the \texttt{async\_slam\_toolbox\_node} provided by the
ROS~2 \texttt{slam\_toolbox} package.

To understand this process, it is first necessary to examine
the main reference frames used during navigation.

\subsubsection{Reference Frames}

Figure~\ref{fig:mapping_frame_tree} illustrates the reference
frame hierarchy.

\begin{figure}[H]
    \centering
    \includegraphics[width=0.8\textwidth]{blank.png}
    \caption{Main reference frames used during mapping:
    \texttt{map}, \texttt{odom}, and \texttt{base\_link}.}
    \label{fig:mapping_frame_tree}
\end{figure}

The hierarchy starts from the fixed \texttt{map} frame, to which
the global costmap is anchored. This frame is connected to the
\texttt{odom} frame through a transform estimated by SLAM.
The \texttt{odom} frame provides the reference for the locally
continuous odometry estimate and is connected to the robot's
\texttt{base\_link} frame.

One of the main tasks of SLAM is to estimate the transform
between \texttt{map} and \texttt{odom}, compensating for
accumulated odometry drift.

\subsubsection{Trajectory and Map Correction}

The SLAM system maintains a graph of robot poses associated with
laser scans. These poses are subsequently optimized to improve
the alignment of the observed LiDAR features.

Consequently, the correction process affects both the
\texttt{map}-to-\texttt{odom} transform and the estimated poses
associated with previously acquired scans. The occupancy map
can therefore be updated using the corrected trajectory,
reducing duplicated or blurred obstacle boundaries.

In principle, a similar correction strategy could also be
applied to previously recorded visual observations, which could
provide additional information for localization. However,
these extensions were not evaluated in the proposed pipeline,
as the results obtained were already sufficiently accurate
and the additional processing would increase computational
and memory requirements.

Nevertheless, extending the mapping framework in this direction
represents a possible avenue for future work.

\noindent\textbf{Reference:}
Macenski, S., and Jambrecic, I. (2021).
\textit{SLAM Toolbox: SLAM for the dynamic world}.
Journal of Open Source Software, 6(61), 2783.
\url{https://doi.org/10.21105/joss.02783}

\subsection{Data Fusion and Outputs: \texttt{semantic\_mapper.py}}
\label{subsec:semantic_mapper}

The \texttt{semantic\_mapper.py} node combines the LiDAR
observations with the visual cost map. LiDAR-detected obstacles
are assigned the maximum cost of 100 and take priority over
visual costs. Laser scan information is also used to populate
regions that have not been observed by the vision pipeline.

The procedure generates three maps:
\begin{enumerate}
    \item \texttt{limo\_map\_complete.pgm}: a complete map
    containing both physical and visual obstacles, with their
    corresponding traversal costs.

    \item \texttt{limo\_map\_laser.pgm}: a map containing only
    LiDAR-detected obstacles, represented with the maximum
    cost of 100.

    \item \texttt{limo\_cv\_obstacle.pgm}: a map containing
    yellow lines and visually detected sidewalks, both
    represented with the maximum cost of 100.
\end{enumerate}

The complete map is used for navigation, while the other two
maps are used for robot localization.

\begin{figure}[H]
    \centering
    \includegraphics[width=\textwidth]{blank.png}
    \caption{Maps generated by \texttt{semantic\_mapper.py}:
    complete cost map, LiDAR-only obstacle map, and visual
    obstacle map.}
    \label{fig:offline_mapping_outputs}
\end{figure}


\section{Autonomous Localization}
\label{sec:autonomous_localization}

\subsection{LiDAR-Based Localization and Its Limitations}
\label{subsec:lidar_localization}

One of the first challenges in autonomous navigation is to estimate
the robot's pose within a given map. One of the most widely used
algorithms for this purpose is Adaptive Monte Carlo Localization
(AMCL), which is implemented in Nav2, the navigation stack for ROS~2.

The underlying principle is to compare LiDAR measurements with a
previously constructed map to estimate a probability distribution
over the robot's possible poses. Ideally, this distribution should
be concentrated around the robot's actual pose, with sufficiently
small uncertainty.

However, in the environments considered here, LiDAR measurements
provide relatively little distinctive information and may
occasionally be unreliable. Both the simulated and real circuits
are fully or partially enclosed by walls forming rectangular
boundaries, whose geometric symmetries can produce ambiguous pose
hypotheses. Moreover, the robot may detect obstacles whose positions
change between experiments or during navigation, introducing
further uncertainty.

To evaluate conventional AMCL under these conditions, a simulation
experiment was conducted in which the robot navigated through one
of the circuits using only LiDAR-based localization. The estimated
covariance remained relatively large throughout the experiment,
as illustrated in Figure~\ref{fig:lidar_amcl_covariance}.

% TODO: Insert the measured covariance threshold and specify
% the corresponding position or orientation component.

\begin{figure}[H]
    \centering
    \includegraphics[width=\textwidth]{blank.png}
    \caption{Estimated pose covariance during navigation using
    conventional LiDAR-based AMCL.}
    \label{fig:lidar_amcl_covariance}
\end{figure}

An important objective of localization is to estimate the robot's
position relative to surrounding obstacles. However, as shown in
Figure~\ref{fig:lidar_amcl_alignment}, the estimated pose did not
always provide sufficiently accurate alignment between sensor
observations and the reference map. Such discrepancies can
negatively affect path planning and obstacle avoidance.

\begin{figure}[H]
    \centering
    \includegraphics[width=\textwidth]{blank.png}
    \caption{Alignment between sensor observations and the
    reference map using conventional LiDAR-based AMCL.}
    \label{fig:lidar_amcl_alignment}
\end{figure}

To address these limitations, information extracted by the computer
vision pipeline is incorporated into localization. Combining LiDAR
measurements with semantic environmental features provides
additional spatial information that can help resolve geometric
ambiguities.

\subsection{Sensor-Fused AMCL}
\label{subsec:sensor_fused_amcl}

The proposed approach extends AMCL to integrate visual information
alongside laser measurements. It preserves the particle-filter
structure while introducing an additional measurement model based
on semantic information extracted from camera images.

Since visual observations are not informative in every situation,
a dedicated assessment mechanism determines whether their
contribution should be included in the particle weight update.

\subsubsection{Particle Representation and Motion Model}

The robot pose distribution is represented by $N$ weighted
particles. Each particle $i$ has a pose
$\boldsymbol{\xi}_{t,i}=(x_{t,i},y_{t,i},\theta_{t,i})$
and a non-negative weight $w_{t,i}$ satisfying
\begin{equation}
    w_{t,i}\geq 0,
    \qquad
    \sum_{i=1}^{N}w_{t,i}=1.
\end{equation}

Each particle represents a possible robot pose in the map frame,
and its weight expresses the relative plausibility of that
hypothesis. The initial particle cloud is generated by applying
Gaussian perturbations around the initial pose estimate.

During navigation, the particles are propagated using the
odometry motion model, with random perturbations accounting
for motion uncertainty:
\begin{equation}
    \boldsymbol{\xi}_{t,i}^{-}
    =
    f\left(
        \boldsymbol{\xi}_{t-1,i},
        \Delta\boldsymbol{\xi}_{\mathrm{odom},t}
    \right)
    +
    \boldsymbol{\varepsilon}_{t,i}.
    \label{eq:amcl_motion_prediction}
\end{equation}

Here, $f(\cdot)$ denotes the deterministic odometry model,
$\boldsymbol{\varepsilon}_{t,i}$ represents motion uncertainty,
and the superscript $-$ denotes the predicted state before
the measurement update.

Once the particles have been propagated, sensor measurements
are used to update their weights. For readability, the time
index $t$ is omitted in the following equations.

\subsubsection{LiDAR Measurement Model}

The LiDAR update follows a likelihood-field approach. For each
particle, valid laser-ray endpoints are projected into the map
frame according to the corresponding pose hypothesis.

For a ray $k$ with measured range $r_k$ and relative bearing
$\alpha_k$, the endpoint coordinates are
\begin{align}
    e_{ik,x}
    &=
    x_{L,i}
    +
    r_k\cos\left(\theta_{L,i}+\alpha_k\right), \\
    e_{ik,y}
    &=
    y_{L,i}
    +
    r_k\sin\left(\theta_{L,i}+\alpha_k\right),
\end{align}
where $(x_{L,i},y_{L,i},\theta_{L,i})$ is the LiDAR pose in the
map frame, obtained from the particle hypothesis and the
sensor extrinsic calibration.

A distance field is computed from the reference occupancy map.
Each cell stores the Euclidean distance to the nearest occupied
obstacle. Endpoints outside the map boundaries are assigned
the maximum cutoff distance.

Let $d_{ik}$ denote the distance associated with the projected
endpoint. Its measurement score is
\begin{equation}
    s_{ik}
    =
    z_{\mathrm{hit}}
    \exp\left(
        -\frac{d_{ik}^{2}}{2\sigma_{\mathrm{hit}}^{2}}
    \right)
    +
    \frac{z_{\mathrm{rand}}}{r_{\max}},
    \label{eq:lidar_measurement_score}
\end{equation}
where $z_{\mathrm{hit}}$ weights obstacle proximity,
$\sigma_{\mathrm{hit}}$ controls the spatial decay,
$z_{\mathrm{rand}}$ weights the uniform random component,
and $r_{\max}$ is the maximum sensor range.

The parameters used are
\begin{equation}
    z_{\mathrm{hit}}=0.5,
    \qquad
    z_{\mathrm{rand}}=0.5,
    \qquad
    \sigma_{\mathrm{hit}}=0.2\,\mathrm{m}.
\end{equation}

The individual ray scores are aggregated into a LiDAR score:
\begin{equation}
    L_i
    =
    1+\sum_{k\in\mathcal{B}}s_{ik}^{3},
    \label{eq:lidar_particle_score}
\end{equation}
where $\mathcal{B}$ is the set of valid rays. The cubic exponent
emphasizes measurements that agree strongly with the map,
while the constant prevents the score from vanishing.

The predicted weight is updated according to
\begin{equation}
    \widetilde{w}_i^{L}
    =
    w_i^{-}L_i.
\end{equation}

Normalization yields
\begin{equation}
    w_i^{L}
    =
    \frac{w_i^{-}L_i}
    {\displaystyle\sum_{j=1}^{N}w_j^{-}L_j}.
    \label{eq:lidar_weight_update}
\end{equation}

\subsubsection{Visual Measurement Model}

The visual measurement model uses a previously generated semantic
map $M_V$ and a point cloud of visual observations acquired online.

The semantic map assigns each cell $c$ one of the following values:
\begin{equation}
    M_V(c)
    =
    \begin{cases}
        0,   & \text{traversable road}, \\
        100, & \text{visual obstacle}, \\
        -1,  & \text{unknown region}.
    \end{cases}
\end{equation}

The visual point cloud is subsampled to reduce computational
overhead. Each retained point $j$ is assigned a binary label:
\begin{equation}
    o_j
    =
    \begin{cases}
        0, & \text{road point}, \\
        1, & \text{obstacle point}.
    \end{cases}
\end{equation}

Temporal compensation accounts for the timestamp offset between
LiDAR and camera observations, aligning the visual point cloud
with the robot pose used for the update.

For particle $i$, a compensated local point $(u_j,v_j)$ is
projected into the map frame as
\begin{align}
    q_{ij,x}
    &=
    x_i^{-}
    +
    u_j\cos\theta_i^{-}
    -
    v_j\sin\theta_i^{-}, \\
    q_{ij,y}
    &=
    y_i^{-}
    +
    u_j\sin\theta_i^{-}
    +
    v_j\cos\theta_i^{-}.
\end{align}

Semantic consistency is evaluated by comparing the observed label
with the map value at the projected position
$\mathbf{q}_{ij}=(q_{ij,x},q_{ij,y})$.
The binary disagreement term is
\begin{equation}
    \delta_{ij}
    =
    o_j\,
    \mathbf{1}\!\left[M_V(\mathbf{q}_{ij})\neq 100\right]
    +
    (1-o_j)\,
    \mathbf{1}\!\left[M_V(\mathbf{q}_{ij})\neq 0\right],
    \label{eq:visual_disagreement}
\end{equation}
where $\mathbf{1}[\cdot]$ is the indicator function.

The overall disagreement score for particle $i$ is
\begin{equation}
    D_i
    =
    \frac{1}{M}
    \sum_{j=1}^{M}\delta_{ij},
    \label{eq:visual_disagreement_score}
\end{equation}
where $M$ is the number of valid visual points used in the
comparison. Consequently, $D_i\in[0,1]$, with lower values
indicating better agreement with the semantic map.

\subsubsection{Visual Information Reliability Assessment}

Visual observations can be ambiguous or contain insufficient
semantic structure. A reliability filter therefore determines
whether the visual update should be applied.

First, the disagreement scores are converted into a normalized
distribution over the particles:
\begin{equation}
    q_i
    =
    \frac{
        \exp\left[-\lambda_V\gamma(D_i-D_{\min})\right]
    }{
        \displaystyle\sum_{j=1}^{N}
        \exp\left[-\lambda_V\gamma(D_j-D_{\min})\right]
    },
    \label{eq:visual_particle_distribution}
\end{equation}
where $D_{\min}=\min_i D_i$, $\lambda_V$ is the visual weight
factor, and $\gamma$ is a sensitivity gain.

The informativeness of the visual observations is evaluated
using the Kullback--Leibler divergence from a uniform
distribution:
\begin{equation}
    I
    =
    \sum_{i=1}^{N}
    q_i\ln\left(\frac{q_i}{1/N}\right).
    \label{eq:visual_kl_divergence}
\end{equation}

A value close to zero indicates that the visual observations
provide little discrimination between the particle hypotheses.

To evaluate spatial dispersion, let
$\mathbf{p}_i=(x_i^{-},y_i^{-})^{\mathsf{T}}$.
The weighted mean position and covariance are
\begin{align}
    \boldsymbol{\mu}_q
    &=
    \sum_{i=1}^{N}q_i\mathbf{p}_i, \\
    \Sigma_q
    &=
    \sum_{i=1}^{N}
    q_i
    \left(\mathbf{p}_i-\boldsymbol{\mu}_q\right)
    \left(\mathbf{p}_i-\boldsymbol{\mu}_q\right)^{\mathsf{T}}.
\end{align}

The maximum positional spread is measured by
\begin{equation}
    \sigma_{\mathrm{pos}}
    =
    \sqrt{\lambda_{\max}(\Sigma_q)},
    \label{eq:visual_position_spread}
\end{equation}
where $\lambda_{\max}(\Sigma_q)$ is the largest eigenvalue
of the covariance matrix.

Angular dispersion is evaluated using circular statistics.
The mean resultant length is
\begin{equation}
    \rho
    =
    \left|
        \sum_{i=1}^{N}q_i
        \exp\left(\mathrm{j}\theta_i^{-}\right)
    \right|,
\end{equation}
where $\mathrm{j}$ is the imaginary unit. The corresponding
circular standard deviation is
\begin{equation}
    \sigma_{\mathrm{yaw}}
    =
    \sqrt{-2\ln\rho}.
    \label{eq:visual_yaw_spread}
\end{equation}

The visual update is accepted only when all the following
conditions are satisfied:
\begin{equation}
    \begin{aligned}
        I &\geq 0.02, \\
        \sigma_{\mathrm{pos}} &\leq 0.5\,\mathrm{m}, \\
        \sigma_{\mathrm{yaw}} &\leq 0.5\,\mathrm{rad}.
    \end{aligned}
    \label{eq:visual_reliability_conditions}
\end{equation}

At least 20 observed obstacle points are also required.
If this requirement or any reliability condition is not met,
the visual update is discarded.

\subsubsection{Sensor Fusion and Particle Weight Update}

When the reliability conditions are satisfied, the LiDAR-updated
weights and visual scores are combined as follows:
\begin{equation}
    w_i^{F}
    =
    \frac{
        (w_i^{L})^{\lambda_L}
        \exp\left(-\lambda_V\gamma D_i\right)
    }{
        \displaystyle\sum_{j=1}^{N}
        (w_j^{L})^{\lambda_L}
        \exp\left(-\lambda_V\gamma D_j\right)
    }.
    \label{eq:fused_particle_weights}
\end{equation}

The parameters $\lambda_L$ and $\lambda_V$ control the relative
influence of the LiDAR-updated distribution and the visual
measurement score. Their values for simulation and real-world
experiments are reported in Table~\ref{tab:amcl_fusion_parameters}.

\begin{table}[H]
    \centering
    \begin{tabular}{lcc}
        \toprule
        \textbf{Parameter}
        & \textbf{Simulation}
        & \textbf{Real Environment} \\
        \midrule
        $\lambda_L$ & 2  & 1 \\
        $\lambda_V$ & 1  & 1 \\
        $\gamma$    & 20 & 20 \\
        \bottomrule
    \end{tabular}
    \caption{Parameters used for LiDAR and visual fusion.}
    \label{tab:amcl_fusion_parameters}
\end{table}

When the reliability checks fail, the visual update is bypassed
and the LiDAR-updated weights are retained:
\begin{equation}
    w_i^{F}=w_i^{L}.
\end{equation}

\subsubsection{Resampling and Iterative Localization}

Following the measurement update, particles are resampled according
to their normalized weights, favouring more plausible hypotheses.

As the robot navigates, motion prediction, measurement updates,
and resampling are repeated. The accepted visual observations
complement the geometric information provided by LiDAR,
helping the filter distinguish between otherwise ambiguous
pose hypotheses.

\subsection{Pipeline Validation}
\label{subsec:localization_validation}

The proposed algorithm was evaluated on the same simulation
circuit used to test conventional LiDAR-based AMCL.

The results showed a substantial reduction in estimated
localization uncertainty, with the position covariance
approaching zero during most of the experiment, as illustrated
in Figure~\ref{fig:fused_amcl_covariance}.

\begin{figure}[H]
    \centering
    \includegraphics[width=\textwidth]{blank.png}
    \caption{Estimated position covariance during navigation
    using Sensor-Fused AMCL.}
    \label{fig:fused_amcl_covariance}
\end{figure}

Furthermore, the visual point cloud showed good spatial alignment
with the corresponding obstacles in the global costmap,
as illustrated in Figure~\ref{fig:fused_amcl_alignment}.

\begin{figure}[H]
    \centering
    \includegraphics[width=\textwidth]{blank.png}
    \caption{Alignment between the visual point cloud and the
    global costmap using Sensor-Fused AMCL.}
    \label{fig:fused_amcl_alignment}
\end{figure}

\section{Path Planning}
\label{sec:path_planning}

Once a global map is available and the robot's pose is estimated
with sufficient precision, a reference path must be generated
to connect the current pose to a desired goal position and
orientation.

The controller computes commands by minimizing a cost function
that accounts for several objectives, including proximity to
the reference path. Consequently, the reference acts as a guide
rather than a trajectory that must be followed exactly.
Strict adherence to a path generated from an imperfect map
may be undesirable, particularly when local observations
indicate that adjustments are necessary.

The global planner must nevertheless account for the robot's
kinematic constraints. This is particularly important in the
experimental road network, where narrow lanes limit the space
available for manoeuvring.

For these reasons, an A*-based planning algorithm is used to
generate paths compatible with the robot's Ackermann steering
geometry.

\subsection{Custom Costmap Inflation}
\label{subsec:custom_inflation}

The planner searches for a path to the goal by minimizing a cost
that depends partly on the global costmap. However, the semantic
map assigns discrete costs to different classes without explicitly
representing preferred positions within a lane. For example,
road cells have zero cost both near the centre of a lane and
close to a road boundary.

This does not necessarily cause the planner to select paths
dangerously close to obstacles. Nevertheless, the lack of spatial
guidance can increase the search effort and produce paths that
do not follow the intended lane structure.

A custom inflation layer is therefore introduced to shape the
cost distribution around obstacles and road boundaries.

\subsubsection{Design Objectives}

Conventional inflation encourages the robot to maintain clearance
from obstacles by assigning larger costs to nearby cells.
For road navigation, obstacle clearance alone is insufficient:
the planner should also favour appropriate lane positioning
while leaving sufficient space for other adjacent vehicles when possible, even where
no centreline is present.

The proposed inflation profile therefore combines obstacle
repulsion with preferred distance bands. These bands guide the
planner towards suitable lateral positions within the road.

\subsubsection{Inflation Cost Profile}

Each non-road known cell generates a distance-dependent cost
profile. The profiles used in simulation and in the real
environment are illustrated in
Figure~\ref{fig:inflation_cost_profiles}.

\begin{figure}[H]
    \centering
    \includegraphics[width=\textwidth]{blank.png}
    \caption{Custom inflation cost profiles used in simulation
    and in the real environment. Regions A--F identify the
    obstacle, repulsion zone, preferred distance bands,
    separating cost peak, and outer cost region.}
    \label{fig:inflation_cost_profiles}
\end{figure}

The profile is divided into six regions:

\begin{description}
    \item[Region A --- Source obstacle.]
    The obstacle cells from which the inflation profile
    is generated.

    \item[Region B --- Obstacle repulsion.]
    A region in which the added cost gradually decreases
    with distance from the obstacle, discouraging paths
    that pass too close to it.

    \item[Region C --- First preferred band.]
    A distance band with zero added inflation cost,
    corresponding to the preferred position in the lane
    closest to the source boundary.

    \item[Region D --- Separating cost peak.]
    A local cost maximum that separates the two preferred
    bands. It introduces a virtual lane separation where
    no explicit centreline is available.

    \item[Region E --- Second preferred band.]
    A second local minimum corresponding to the other lane.
    This band is particularly useful in wider junctions
    with sparse road markings, where it encourages more
    orderly turning manoeuvres.

    \item[Region F --- Outer cost region.]
    Beyond the second preferred band, the added cost
    increases gradually until reaching a saturation value,
    where configured. This discourages unnecessary
    deviations from the preferred road region.
\end{description}

Where multiple inflation profiles overlap, the contribution from the nearest source cell is retained. The resulting inflation cost is
combined with the underlying semantic costmap by taking
the maximum of the two values:
\begin{equation}
    C_{\mathrm{final}}(c)
    =
    \max\left(
        C_{\mathrm{semantic}}(c),
        C_{\mathrm{inflation}}(c)
    \right),
    \label{eq:inflation_cost_combination}
\end{equation}
where $c$ denotes a map cell. Consequently, a zero inflation
cost does not remove an existing semantic cost.

Figure~\ref{fig:inflated_global_costmap} shows the resulting
cost distribution over the global map, with selected examples
of the six regions highlighted.

\begin{figure}[H]
    \centering
    \includegraphics[width=\textwidth]{blank.png}
    \caption{Global costmap after custom inflation.
    Annotated areas illustrate examples of regions A--F
    within the road network.}
    \label{fig:inflated_global_costmap}
\end{figure}

\subsection{The Algorithm: \texttt{rviz\_goal\_bridge.py}}

The main planning pipeline is illustrated in
Figure~\ref{fig:goal_bridge_flowchart}.

\begin{figure}[htbp]
    \centering
    \includegraphics[width=0.95\textwidth]{blank.png}
    \caption{Flowchart of the start--goal adjustment and path-planning procedure.}
    \label{fig:goal_bridge_flowchart}
\end{figure}

The user selects a target pose in RViz2, which is received by the
\texttt{rviz\_goal\_bridge.py} node. Although this request could be
forwarded directly to the path-planning service, small inaccuracies
in either the estimated starting pose or the selected goal may
prevent the planner from finding a feasible solution. For instance,
the robot footprint may slightly overlap an obstacle at either
endpoint, or the requested poses may be difficult to connect under
the robot's kinematic constraints.

To improve robustness, the bridge introduces a
\textbf{start--goal adjustment procedure}. This procedure generates
alternative poses around both endpoints, checks them for collisions,
and submits candidate start--goal pairs to the planner, prioritising
small deviations from the original request.

\subsubsection{Candidate Pose Generation}

Candidate positions are sampled on concentric rings centred on each
reference position. The radius of each ring is an integer multiple
of the radial sampling interval, $\Delta r$, up to a maximum search
radius, $R$.

Positions are distributed uniformly along each ring, with their
number increasing with its circumference to maintain an
approximately constant sampling distance. Each ring contains at
least eight positions. The original reference position is also
included.

At each sampled position, the algorithm evaluates 16 orientations,
spaced at $22.5^\circ$ intervals relative to the reference heading.
These include the original orientation and offsets covering a full
revolution.

With a search radius of $0.75\,\mathrm{m}$ and a radial interval of
$0.10\,\mathrm{m}$, the outermost sampled ring lies at
$0.70\,\mathrm{m}$. This produces 181 positions and 16 orientations
per position, giving a total of \textbf{2,896 candidate poses for
each endpoint}, before collision filtering.

\subsubsection{Candidate Ranking and Collision Filtering}

Each candidate receives a score combining its translational
displacement and its absolute angular deviation from the reference
pose. Candidates are then sorted in ascending order of this score,
favouring poses close to the original request.

Collision checks use a rectangular robot footprint measuring
$0.322\,\mathrm{m} \times 0.220\,\mathrm{m}$. The footprint width
includes a safety margin of $2.5\,\mathrm{mm}$ on each side relative
to the physical robot. A candidate is retained only if its footprint
does not overlap either a LiDAR obstacle or a boardwalk obstacle.

The procedure therefore produces two ordered lists of collision-free
candidates: one for the start and one for the goal.

\subsubsection{Start--Goal Pair Selection}

The bridge combines the two lists and tests candidate pairs in a
diagonal sequence. It first considers the highest-ranked start and
goal, followed by pairs with progressively larger sums of their
list indices:
%
\[
\begin{aligned}
(s_0,g_0)
&\rightarrow (s_0,g_1)
\rightarrow (s_1,g_0) \\
&\rightarrow (s_0,g_2)
\rightarrow (s_1,g_1)
\rightarrow (s_2,g_0)
\rightarrow \cdots
\end{aligned}
\]

This ordering prioritises pairs drawn from the beginning of both
lists, allowing adjustments to either endpoint before considering
more distant alternatives.

\subsubsection{Path Planning}

Each candidate pair is submitted to \textbf{SMAC Hybrid A*},
provided by the Navigation2 \texttt{smac\_planner} package, using
the \textbf{Dubins motion model}.

Hybrid A* follows the search principle of conventional A*, but
expands states using motion primitives that account for the
vehicle's heading and turning constraints. Collision checking also
considers the robot footprint along the candidate path, rejecting
configurations that intersect obstacles classified as
non-traversable.

The Dubins motion model restricts the search to \textbf{forward
motion}, so the resulting global paths do not include reversing
manoeuvres.

The bridge tests successive start--goal pairs until the planner
finds a feasible path or reaches the maximum number of attempts.
This limit is set to \textbf{48} in both the simulation and
real-robot configurations.

Figure~\ref{fig:goal_bridge_adjustment} illustrates the benefit of
endpoint adjustment. The turquoise path is successfully generated
between nearby valid poses, even though the original starting pose
and the user-selected goal would prevent a direct planning request
from succeeding.

\begin{figure}[H]
    \centering
    \includegraphics[width=0.95\textwidth]{blank.png}
    \caption{Successful path generation through start--goal adjustment.
    The turquoise path connects nearby valid poses despite the
    infeasibility of the original endpoints.}
    \label{fig:goal_bridge_adjustment}
\end{figure}

(source for smac:\url{https://arxiv.org/html/2401.13078v2})

\section{Online Local Mapping}
\label{sec:online_mapping}

The controller requires information about both the robot's position
relative to the reference trajectory and the physical and visual
obstacles detected during navigation.

However, simply combining the latest semantic point cloud with
the current LiDAR observations is insufficient. As illustrated
in Figure~\ref{fig:visual_blind_region}, the camera detects
obstacles ahead of the robot but leaves a blind region immediately
around its footprint. Visual obstacles therefore disappear from
the incoming observations as the robot approaches them, even
though they remain relevant to navigation. A short-term spatial
memory is consequently introduced to retain selected observations
after they leave the camera's field of view.

\begin{figure}[H]
    \centering
    \includegraphics[width=\textwidth]{blank.png}
    \caption{Visual observations produced by the perception
    pipeline, showing the blind region immediately around
    the robot's footprint.}
    \label{fig:visual_blind_region}
\end{figure}

The area ahead of and around the robot is divided into three
regions, all defined in the \texttt{base\_link} frame:

\begin{enumerate}
    \item \textbf{Local map rectangle (green):}
    defines the full spatial extent of the local visual map.

    \item \textbf{Visual trapezoid (yellow):}
    approximates the ground region currently visible
    to the camera.

    \item \textbf{Inner trapezoid (red):}
    is obtained by offsetting selected boundaries of the
    visual trapezoid inwards and encloses the interior
    of the observed region.
\end{enumerate}

\begin{figure}[H]
    \centering
    \includegraphics[width=\textwidth]{blank.png}
    \caption{Regions used for local visual mapping:
    local map rectangle (green), visual trapezoid (yellow),
    and inner trapezoid (red).}
    \label{fig:local_mapping_regions}
\end{figure}

\subsection{The Algorithm: \texttt{local\_ctrl\_map.py}}
\label{subsec:local_ctrl_map}

The algorithm processes the following semantic classes:

\begin{enumerate}
    \item \texttt{LABEL\_YELLOW}
    \item \texttt{LABEL\_BOARDWALK}
    \item \texttt{LABEL\_INTERIOR\_BOARDWALK}
\end{enumerate}

These classes represent the most relevant visual features for
local navigation. Boardwalk regions constitute the main visual
obstacles, while yellow markings provide distinctive guidance
for lane following.

Let:
\begin{itemize}
    \item $A_c(t)$ denote the set of visual points belonging
    to class $c$ observed at time $t$;

    \item $M(t)$ denote the set of stored points at time $t$;

    \item $G$, $Y$, and $R$ denote the green rectangle,
    yellow trapezoid, and red trapezoid, respectively;

    \item $L(t)$ denote the output local map.
\end{itemize}

For simplicity, the following description uses discrete update
indices $t$. The implementation does not operate at fixed time
intervals: the elapsed time $\Delta t$ between consecutive updates
is computed from their timestamps.

The procedure uses the \texttt{base\_link} and \texttt{odom}
reference frames. Incoming visual observations are expressed
in \texttt{base\_link}:
\begin{equation}
    A(t)
    =
    \left\{
        \left(x_i(t),y_i(t)\right)^{\mathrm{base\_link}}
        \;\middle|\;
        i\in\mathcal{I}(t)
    \right\},
    \label{eq:local_mapping_observations}
\end{equation}
where $\mathcal{I}(t)$ indexes the points produced by the visual
pipeline at update $t$.

At each update, the following operations are performed.

\subsubsection{Observation Selection}

A subset $B(t)$ is selected from the incoming observations.
Only points belonging to classes 1 and 3 and located between
the visual and inner trapezoids are retained:
\begin{equation}
    B(t)
    =
    \left(
        \bigcup_{c\in\{1,3\}}A_c(t)
    \right)
    \cap
    \left(Y\setminus R\right).
    \label{eq:local_mapping_selection}
\end{equation}

Here, indices 1 and 3 refer to \texttt{LABEL\_YELLOW} and
\texttt{LABEL\_INTERIOR\_BOARDWALK}, respectively.

\subsubsection{Memory Update}

The selected points are transformed into \texttt{odom}
to facilitate subsequent reprojection. They are then added
to memory:
\begin{equation}
    M(t)
    =
    M(t^{-})
    \cup
    B^{\mathrm{odom}}(t),
    \label{eq:local_mapping_memory_update}
\end{equation}
where $M(t^{-})$ denotes the memory immediately before the
current update. The observations in $B^{\mathrm{odom}}(t)$
are augmented with confidence values when inserted into memory.

Each stored point is associated with a confidence value:
\begin{equation}
    M(t)
    =
    \left\{
        \left(x_i,y_i,c_i(t)\right)^{\mathrm{odom}}
        \;\middle|\;
        i\in\mathcal{I}_M(t),
        \quad c_i(t)\in[0,1]
    \right\},
    \label{eq:local_mapping_memory_representation}
\end{equation}
where $\mathcal{I}_M(t)$ indexes the stored points.
The semantic class of each point is also retained.

\subsubsection{Spatial Downsampling}

The stored points are spatially downsampled using a grid
covering their spatial extent. At most one point per semantic
class is retained in each cell, prioritising the most recent
observations.

This reduces redundant information while preserving the
semantic structure of the local environment.

\subsubsection{Memory Size Limitation}

An upper limit is imposed on the number of stored points
to bound processing time and memory consumption. If this
limit is exceeded, the points with the lowest confidence
are removed.

\subsubsection{Transformation into the Robot Frame}

The stored points are transformed from \texttt{odom} back
into the current \texttt{base\_link} frame. Their positions
can then be evaluated relative to the robot and the local
mapping regions.

\subsubsection{Confidence Decay}

The confidence of each stored point decreases exponentially
according to
\begin{equation}
    c_i(t+1)
    =
    c_i(t)
    \exp\left(
        -\lambda\,m(t)\,g_i(t)\,\Delta t
    \right),
    \label{eq:local_mapping_confidence_decay}
\end{equation}
where:
\begin{itemize}
    \item $\lambda=0.10\,\mathrm{s}^{-1}$ is the base decay rate;

    \item $m(t)\in[0,1]$ is a velocity-dependent factor;

    \item $g_i(t)$ is a spatial factor;

    \item $\Delta t$ is the elapsed time between consecutive
    updates, computed from their timestamps.
\end{itemize}

The velocity-dependent factor reduces the decay when the robot
moves slowly, allowing stored observations to cover the nearby
blind region before their confidence becomes too low.

The spatial factor is defined as
\begin{equation}
    g_i(t)
    =
    \begin{cases}
        3, & \text{if the point lies inside }Y, \\
        1, & \text{otherwise}.
    \end{cases}
    \label{eq:local_mapping_spatial_decay}
\end{equation}

Observations therefore decay faster within the currently visible
region, where fresh measurements are available.

\subsubsection{Spatial and Confidence Filtering}

A further filtering step removes points outside the local map,
inside the inner trapezoid, or below the confidence threshold.
The retained set is
\begin{equation}
    \begin{aligned}
        \widetilde{M}(t)
        =
        \bigl\{
            (x,y,c)
            \;\big|\;&
            (x,y,c)\in M^{\mathrm{base\_link}}(t), \\
            &(x,y)\in G\setminus R,
            \quad c>0.3
        \bigr\}.
    \end{aligned}
    \label{eq:local_mapping_filtered_memory}
\end{equation}

This preserves sufficiently reliable observations in regions
where they can supplement the current camera data.

\subsubsection{Local Map Publication}

The output local map combines the current observations
with the filtered memory points:
\begin{equation}
    L(t)
    =
    A(t)\cup\widetilde{M}(t).
    \label{eq:local_mapping_output}
\end{equation}

In this expression, the union combines the spatial observations
and their semantic labels; confidence values are used internally
to filter the stored points.

The resulting local map is then rasterised and published as
a \texttt{nav\_msgs/msg/OccupancyGrid} message for use by
the controller.

Spatial and temporal filtering are applied at several stages
to keep the computational requirements bounded. These operations
must nevertheless retain enough points to represent nearby
obstacles and road markings adequately, particularly within
the camera's blind region.

\begin{figure}[H]
    \centering
    \includegraphics[width=\textwidth]{blank.png}
    \caption{Local costmap generated by
    \texttt{local\_ctrl\_map.py}, combining current visual
    observations with retained memory points.}
    \label{fig:local_mapping_result}
\end{figure}

% TODO: Report the average number of retained points,
% processing times, and publishing frequency.

\section{Local Planning and Control}
\label{sec:local_planning_control}

The local controller follows the geometric path produced by the global
planner while adapting to obstacles detected during operation. At each
control cycle, it uses the current robot state, a sequence of reference
poses and the latest local costmap to select a command that respects
the vehicle's kinematic and dynamic limits.

\subsection{Local Map for Control}
\label{subsec:local_map_for_control}

The controller uses the local costmap described in
Section~\ref{sec:online_mapping}. The map combines current visual
observations with retained observations, which preserve information
about nearby road markings and obstacles when they move into the
camera's blind region. During trajectory evaluation, the costmap
provides the spatial information used to check candidate motions and
account for nearby obstacles.

\subsection{Controller Design}
\label{subsec:controller_design}

The control problem is formulated as a sampling-based nonlinear model
predictive control (MPC) problem. At each cycle, the controller
predicts how different sequences of speed and steering commands will
move the vehicle over a finite horizon. It then applies the first
command of the selected sequence and repeats the procedure using
updated state and map information.

Nav2 for ROS~2 Humble provides a Model Predictive Path Integral (MPPI)
controller through the \texttt{nav2\_mppi\_controller} package.
At each control cycle, MPPI samples perturbations of a nominal command
sequence, rolls out the resulting trajectories using a vehicle model,
and evaluates them with cost functions. It then updates the nominal
sequence by combining the sampled perturbations with cost-dependent
weights and applies the first command of the updated sequence.\footnote{%
    \href{https://api.nav2.org/nav2-humble/html/md_nav2_mppi_controller_README.html}
    {Nav2 Humble documentation for the MPPI controller}.}

MPPI was tested on a ROS~2 Humble branch of the project and showed
satisfactory behaviour in the initial tests. The onboard software
uses ROS~2 Foxy, whose Nav2 distribution does not provide this
controller. The implementation therefore builds on Nav2's DWB
controller from the \texttt{nav2\_dwb\_controller} package. DWB
provides integration with Nav2 and the local costmap, along with
plugin interfaces for trajectory generation and scoring. Its standard
trajectory generators support differential-drive and omnidirectional
robots, and the controller returns a single velocity command at each
cycle.\footnote{%
    \href{https://api.nav2.org/nav2-humble/html/md_nav2_dwb_controller_README.html}
    {Nav2 Humble documentation for the DWB controller}.}

Custom C++ plugins adapt this framework to the LIMO's Ackermann
steering. At each control cycle, the resulting controller
\begin{enumerate}
    \item generates candidate command sequences and simulates the
          corresponding trajectories from the current robot state,
          subject to the vehicle's kinematic and dynamic limits,
    \item rejects trajectories that violate environmental constraints
          represented in the local costmap,
    \item scores the remaining trajectories using the configured cost
          functions, and
    \item applies the first command of the lowest-cost candidate.
\end{enumerate}

Selecting a single candidate keeps the search tractable on the onboard
computer and makes it possible to inspect the sequence selected at
each cycle during debugging.

\subsection{Trajectory Generation}
\label{subsec:control_trajectory_generation}

The global path specifies reference positions and headings without
assigning a speed to each pose. Even along a straight segment, the
appropriate speed depends on the robot's current state and the
obstacles ahead. The controller may need to decelerate, come to a stop,
or temporarily deviate from the path to pass an unexpected obstacle.
In the bicycle model, the yaw rate depends on both the forward speed
and the steering angle. The timing and order of the commands over the
prediction horizon therefore affect the resulting vehicle pose. The
candidate set must cover a range of manoeuvres while keeping the
number of rollouts manageable within each control cycle.

\subsubsection{Basic Candidates}

The generator considers three longitudinal actions: decelerate,
maintain speed, and accelerate. Each is paired with seven steering
options: reduce the current front-wheel steering angle by one of
three sampled increments, hold it, or increase it by one of the same
three increments. Let \(\delta_0\) be the current steering angle.
The sampled increments are computed as
\begin{equation}
    \Delta_i = \delta_{\max}\frac{i}{3}\operatorname{rand}(),
    \qquad i\in\{1,2,3\},
    \label{eq:control_steering_increments}
\end{equation}
where \(\delta_{\max}\) sets the scale of the sampled steering
increments and \(\operatorname{rand}()\) supplies a random scaling
factor.

\begin{table}[H]
    \centering
    \small
    \begin{tabular}{@{}lccc@{}}
        \toprule
        Longitudinal action & Reduce steering & Hold steering & Increase steering \\
        \midrule
        Decelerate     & \(\delta_0-\Delta_i\) & \(\delta_0\) & \(\delta_0+\Delta_i\) \\
        Maintain speed & \(\delta_0-\Delta_i\) & \(\delta_0\) & \(\delta_0+\Delta_i\) \\
        Accelerate     & \(\delta_0-\Delta_i\) & \(\delta_0\) & \(\delta_0+\Delta_i\) \\
        \bottomrule
    \end{tabular}
    \caption{Basic combinations of longitudinal and steering actions.
    In the first and last steering columns, \(i\in\{1,2,3\}\). Each
    longitudinal action is therefore paired with seven steering options.}
    \label{tab:control_basic_actions}
\end{table}

The three longitudinal actions and seven steering options yield
\(3\times 7=21\) basic candidates.

\subsubsection{Steering Sequences for Obstacle Avoidance}

The generator adds 27 candidates that pair the three longitudinal
actions with nine steer--countersteer--steer sequences. Changing the
steering direction within the prediction horizon lets the vehicle
move laterally and then recover a heading compatible with the path.
These candidates are particularly useful when the robot must change
lanes to pass an obstacle, as illustrated in
Figure~\ref{fig:control_countersteering}.

\begin{figure}[H]
    \centering
    \includegraphics[width=0.65\textwidth]{blank.png}
    \caption{Examples of steer--countersteer--steer trajectories that
    produce lateral displacement while recovering the vehicle heading.}
    \label{fig:control_countersteering}
\end{figure}

\subsubsection{Additional Candidates and Trajectory Rollout}

The generator also considers the command sequence selected at the
previous control cycle, shifted forward by one time step, and a
zero-speed sequence. If the robot is moving slowly enough, it adds
three low-speed start sequences to help the vehicle align with the
global path. These sequences are especially useful when the start
pose used by the global planner was adjusted during the
start--goal adjustment procedure described in
Section~\ref{sec:path_planning}.

With the two additional sequences, the candidate set contains
\(21+27+2=50\) sequences. Including the three low-speed starts brings
the total to 53. Each sequence is rolled out from the current robot
state using the bicycle model while enforcing the vehicle's kinematic
and dynamic limits. Figure~\ref{fig:control_trajectory_candidates}
illustrates the resulting candidate trajectories.

\begin{figure}[H]
    \centering
    \includegraphics[width=\textwidth]{blank.png}
    \caption{Example candidate trajectories generated in one control
    cycle. The legend identifies the basic candidates, the
    steer--countersteer--steer candidates and the additional candidates.}
    \label{fig:control_trajectory_candidates}
\end{figure}

\section{Simulation Results}

% ==========================================
% PART II - EXPERIMENTAL IMPLEMENTATION
% ==========================================

\part{Experimental Implementation}

% ==========================================
% CHAPTER 5 - REAL WORLD IMPLEMENTATION
% ==========================================

\chapter{Real-World Implementation}

\section{Experimental Setup}

\section{Computer Vision Pipeline}

\section{Mapping and Localization}

\section{Planning and Control}

\section{Experimental Results}

% ==========================================
% CHAPTER 6 - CONCLUSIONS
% ==========================================

\chapter{Conclusions}

\section{Main Achievements}

\section{Limitations and Future Developments}

\end{document}


